\chapter{Experiment 3: Full-Cohort Waveform Delay-Tensor Analysis}
\label{ch:v3}

\section{Why use the leftovers?}

The controlled cohort is valuable for comparability but excludes signal from 19
additional people and thousands of valid windows. V3 therefore creates a new
analysis index over both stage-02 primary and leftover partitions. It does not
merge or overwrite those partitions. Uniform CAR, QC, and boundary rules yield
16,824 retained windows from all 88 people.

The key imbalance rule changes from equal window counts to equal \emph{person}
influence. For training label $y_i$, the subject weight is

\begin{equation}
\omega_i=\frac{N_{\mathrm{train}}}
{3N_{\mathrm{train},y_i}}.
\label{eq:class-weight}
\end{equation}

Each diagnosis receives one third of total loss mass. Within a person, windows
are averaged, not summed; equivalently, each window contributes its person's
weight divided by that person's retained-window count. There is no duplication,
SMOTE, or synthetic \FTD{} data.

\section{Waveform branch}

The waveform representation is the delay tensor and exact moment method from
Eq.~\ref{eq:moment-energy}. The unsupervised projector first forms a class-
balanced training moment. For TT, a channel factor of rank 6 is obtained from
the channel marginal; a second eigensolve in the reduced channel-lag space
produces 8 or 16 filters. Tucker uses channel and lag ranks $(4,4)$ or $(6,6)$.
Log pooled filter energies become features.

The code numerically checks that moment-derived energy equals explicit raw
delay-tensor contraction and that its TT factors agree with an equivalent
TT-SVD construction on synthetic data.

\section{Supervised TT}

The supervised model parameterizes filters as

\begin{equation}
b_{c\ell k}=\sum_{r=1}^{R_c}U_{cr}V_{r\ell k},
\end{equation}

then maps standardized log energies to three softmax logits. V3 jointly optimizes
$U$, $V$, classifier head $H$, and bias using class-weighted cross-entropy:

\begin{align}
\mathcal{L}={}&\sum_i \tilde{\omega}_i
\left[\log\sum_{g=1}^3 e^{q_{ig}}-q_{i,y_i}\right]
+\frac{\lambda_H}{2}\|H\|_F^2\\
&+\frac{10^{-3}}{2}
\left(\|U-U_0\|_F^2+\|V-V_0\|_F^2\right).
\end{align}

The training-only unsupervised TT initializes $U_0,V_0$ and energy scaling.
L-BFGS-B has a fixed 100-iteration budget; head penalties 0.01 and 0.1 are inner
CV candidates. A reported convergence flag is an optimizer status, not a proof
of global optimality for this nonconvex problem.

\section{Spectral control and candidate set}

The same 418 spectral definitions as V2 are recomputed on all retained V3
windows. Training-only ANOVA selects 16 or 32. Wave TT ranks 8/16, waveform
Tucker ranks 4/6, and the spectral settings are each paired with logistic
regression ($C=0.1,1$) or RBF SVM ($C=1,10$). Together with two supervised TT
penalties, V3 has 26 candidate configurations.

Evaluation uses nested three-repeat five-fold outer / three-fold inner
participant CV, outer seed 9060 and inner seed $90600+$ outer index. The primary
endpoint is again the complete inner-selected pipeline.

\section{Results}

\begin{table}[htbp]
\centering
\begin{tabular}{lrrr}
\toprule
\textbf{V3 family} & \textbf{Mean BA} & \textbf{Accuracy} & \textbf{Macro F1}\\
\midrule
Inner-selected complete pipeline (primary) & 52.0\% & -- & --\\
Spectral control & 55.8\% & -- & --\\
Waveform Tucker & 51.8\% & -- & --\\
Waveform TT & 50.9\% & -- & --\\
Supervised TT & 49.9\% & -- & --\\
\bottomrule
\end{tabular}
\caption{Full-cohort V3 results. The spectral control is the strongest observed family.}
\label{tab:v3-results}
\end{table}

The negative result is informative: more raw-looking waveform structure does
not automatically improve classification. The delay representation captures
short-lag second-order energy but removes absolute phase, local ordering, and
long-range temporal dynamics. The spectral control may better match the stable
disease-related signal available in this small cohort.

\section{Optimization-budget sensitivity}

All 120 supervised TT fits in the original V3 benchmark reach the 100-iteration
limit. A post-hoc sensitivity keeps the already specified head penalty 0.1 and
compares 100 with 2,000 iterations on the same 15 full outer-training splits and
10 smaller learning-curve subsets. Across full splits, 100 iterations averages
55.0\% balanced accuracy and 2,000 averages 53.7\%. Only 3 of 25 long-budget
fits report convergence. More iterations alone therefore do not resolve the
model's weakness; these values are diagnostic and do not replace the V3 primary
endpoint.

The fixed-model learning curve uses approximately 27, 48, and 70 training people
and obtains 42.0\%, 40.9\%, and 55.2\% held-out balanced accuracy. The curve is
non-monotone and based on one CV repetition, so it cannot estimate the number of
additional participants required.

\section{Verification}

V3 checks reconstruct sampled source windows, confirms channel order and class
weight totals, replays saved models and scores, verifies subject-disjoint splits,
preserves every V2 window, excludes all prior artifact IDs, and fingerprints the
unchanged primary and leftover arrays.

